\scanpage{30}
comme un \deleted{\unsure{élément de} $\unsure{\mathcal A_E}$} \inserted{élément de} $\mathcal A_E$. On a alors
\[
(2)\qquad\Delta(1+A)=\Delta_{\deleted{E}}^{\inserted{E}}(1_E+A)
\qquad\text{si}\quad\deleted{A\in\mathcal A_E}\ \inserted{A\in\mathcal A_E}.
\]

\deleted{et si \unsure{A est} (\unsure{que d'ailleurs} \illegible{3} de la \unsure{forme} $\unsure{1+|A+}$ opérateur à trace finie) \unsure{l'intégrale}}

Comme on vérifie aussitôt.

\uline{Prop. 8}
\margin{$\Big|$}
\diagram{\begin{tikzcd}[row sep=1.2em,column sep=4em]
\Delta(A)=\Delta(A^*) & \\
\Delta(AB)=\Delta(A)\Delta(B) & \\
\Delta(\unsure{U|1})=1 \arrow[uu,bend left=55] & \text{si }A,B\in1+\underline a
\end{tikzcd}}

Seule la deuxième formule est non triviale.

Par raison de continuité, on peut supposer qu'il existe $E\in\underline a$ tel que $A\in\mathcal A_E$, $B\in\mathcal A_E$; \unsure{grâce à|puis} (2), ceci nous ramène \deleted{\unsure{au cas}} \inserted{au cas} où \deleted{$\mathcal A$ est à trace finie} \inserted{dans $\mathcal A$ la trace est finie}. \unsure{On}

\deleted{peut supposer de plus (par continuité) $A$ et $B$ \unsure{strict. positifs}.
\[
\deleted{\Delta(A)^2=\Delta(A^*A)}
\]
\deleted{$\Delta(\unsure{A})$}
\margin{$\unsure{u}$}
\[
\begin{aligned}
\Delta(A)^2&=\Delta(AA^*)\\
\Delta(A^*)^2&=\Delta(A^*A)\\
\Delta(AB)^2&=\Delta(B^*\unsure{A^*A|A^2}B)=\Delta(B^*HB)
\overset{?}{=}\Delta(A^*A)\Delta(B^*B).
\end{aligned}
\]
Soit $B^*B=\unsure{B_1^2|L^2}$; posons
\[
\Delta(B^*HB)=\Delta(LHL)
\]
On est ramené à prouver
\[
\Delta(LHL)
\]
}
\mspage{31}
